\scanpage{184}\mspage{40}
% archive sheet number: 184

pour des préfaisceaux à valeurs dans une catégorie $M$ quelconque (sous réserve d'existence des opérations \deleted{\unsure{de quelque nature}} d'image directe \inserted{cohomologique ou homologique} \unsure{envisagées|considérées}).

Cela se voit p. ex. en notant que

\[
W_0\text{-équivalence}\ \Longrightarrow\ \mathbb{D}_M\text{-équivalence}\tag*{(3.10.)}
\]

% slanted
\margin{\unsure{pour être tout|pour être} à fait \unsure{correct}, \inserted{\unsure{pour|parmi}} \deleted{il faut} ``\uline{\unsure{toute}}'' $W_0$-notion, excepter les deux notions de $W_0$-fibration (forte, et faible).}

pour tte \inserted{$\mathcal{U}$-}catégorie $M$, d'où résulte que \unsure{une|toute} \inserted{(*)} ``$W_0$-notion'' implique la $\mathbb{D}_M$-notion correspondante, et permet donc d'appliquer à $\mathbb{D}_M$ les résultats \unsure{le|la} concernant. (Il y a bien sûr un argument \unsure{direct moins|direct mais|d'autant} \unsure{savant|suivant}, consistant à se ramener au cas où $M$ est petite, et à utiliser alors les plongements pl. fidèles $M\hookrightarrow M^\wedge$ resp. $M\hookrightarrow((M^\circ)^\wedge)^\circ$.)

Je rappelle pour mémoire une (\unsure{troisième} ?) caractérisation des flèches $u:\deleted{\unsure{I}}I'\longrightarrow I$ $W_0$-propres (resp $W_0$-lisses) : ce sont celles \deleted{\unsure{p}} pour lesquelles le foncteur chgt. de base $u^*:\mathrm{Cat}_{/I}\longrightarrow\mathrm{Cat}_{/I'}$ transforme flèches $W_0$-coasphériques (resp. $W_0$-asphériques) en itou. Cela gagne en intérêt, du fait que la notion de foncteur $W_0$-coasphérique (resp. $W_0$-asphérique) s'interprète comme celle de foncteur cofinal (resp. co-cofinal), cf plus bas.
